\section{Registered Benchmark and Protocol}
\subsection{Matched IEEE 39-Bus Testbed}

The principal case is a reconciled IEEE 39-bus system with GFL converter
plants at buses 36, 37, and 38. The population remains GFL-only throughout the
campaign; no GFL-to-grid-forming model substitution is part of the action
space. The ANDES implementation uses a symbolic--numeric DAE framework
\cite{cui2021andes}. Its static data were reconciled to a matched ParaEMT
model before the main campaign.

The E02C prerequisite verified the common post-conversion operating point,
network-wide stationarity, converter active/reactive power, current and PLL
channels, adjacent-timestep convergence, and generalized-QZ eigenpair
residuals. All seven registered readiness gates passed after retaining the
canonical 0.97143 transformer tap. This establishes that the benchmark is
suitable for the ANDES experiments below. It does \emph{not} provide EMT
validation of any externality, modal mechanism, or intervention.

The 100-MVA ac solve uses bus 39 as its reference/slack bus. GFL buses 36--38
are fixed-$P/Q$ injections: their legacy PV voltage equations are disabled and
their registered reactive injections are fixed, while the remaining
synchronous PV controls and the slack close the network equations. Thus the
A7 active-power increment is balanced by the reference generator through the
power flow. Every accepted vertex was outside the implemented converter
current and controller-saturation limits. Across the recorded E03 vertices,
minimum bus voltage was 0.84056 pu and maximum GFL current was 0.90842 pu on
the device base, below the 1.2-pu current limit.

\subsection{Finite Engineering Actions}

Table~\ref{tab:actions} lists the eight actions frozen before discovery. The
controller actions preserve model class and state dimension. For PLL
bandwidth actions, the proportional and integral gains follow
$K_p=K_{p0}r$ and $K_i=K_{i0}r^2$ with $r=1+0.2a_i$. The current-command
bandwidth actions use the corresponding registered REGCP interpolation.
A7 changes the GFL37 active-power reference by 20 MW (0.20 on the 100-MVA
system base, 0.02 on the 1000-MVA device base), and A8 increases the series
reactance of branch 25--37 by 10\%. The three GFL ratings are 1000 MVA, with
baseline active injections of 560, 540, and 830 MW. The frozen PLL natural
frequency changes from 2.0 to 2.4 Hz; the current-command time constant changes
from 0.0200 to 0.0167 s; and the branch-25--37 reactance changes from 0.0232 to
0.02552 pu. A8 was
chosen by a static network rule: branch 25--37 is the median-reactance member
of the three transformer branches incident to the GFL cluster. No dynamic
result informed the choice.

\begin{table}[t]
\caption{Frozen Finite Action Set}
\label{tab:actions}
\centering
\footnotesize
\begin{tabular}{clcl}
\toprule
ID & Family & Location & Full action\\
\midrule
A1 & PLL bandwidth & bus 36 & $+20\%$ natural frequency\\
A2 & PLL bandwidth & bus 37 & $+20\%$ natural frequency\\
A3 & PLL bandwidth & bus 38 & $+20\%$ natural frequency\\
A4 & Current-command BW & bus 36 & $+20\%$ bandwidth\\
A5 & Current-command BW & bus 37 & $+20\%$ bandwidth\\
A6 & Current-command BW & bus 38 & $+20\%$ bandwidth\\
A7 & Active-power reference & bus 37 & $540\rightarrow560$ MW\\
A8 & Network reactance & line 25--37 & $0.0232\rightarrow0.02552$ pu\\
\bottomrule
\end{tabular}
\end{table}

The 24 preregistered operating points comprise 16 discovery and eight locked
holdout seeds. Each point evaluates all 28 pairs and 56 triples, for 672 pair
and 1,344 triple cases. A failed ac solution, initialization, descriptor
build, or smooth-path check invalidates the case and remains in the failure
ledger; cases are never silently dropped.

\subsection{Spectral Functional and Numerical Uncertainty}

The primary band is 0.1--30 Hz. Frequency integration uses Gauss--Legendre
quadrature in log frequency with 48 nodes and a 96-node reference. A
conditioning-only pilot selected the first admissible fixed contour shift,
$\sigma=0$, from a preregistered list. A 0.1--50 Hz band is retained only as a
sensitivity analysis. Connected integration uses the frozen action-hypercube
rule, complete pair/triple partitions, and sparse active-column trace solves;
the implementation never forms the full $T^{-1}$ for the campaign.

The direct uncertainty $\epsilon_{\mathrm{direct}}$ is a conservative sum of
case-specific frequency refinement, sparse-log-determinant, and power-flow/
Jacobian terms. A direct coefficient is called resolved only if
\begin{equation}
|\ext{S}\Phi|>5\epsilon_{\mathrm{direct}}.
\label{eq:resolve}
\end{equation}
The reconstruction criterion is
\begin{equation}
|\ext{S}\Phi_{\mathrm{conn}}-\ext{S}\Phi_{\mathrm{direct}}|
\le 5\epsilon_{\mathrm{combined}},
\label{eq:recongate}
\end{equation}
where the combined budget also includes derivative, action-grid, hypercube,
and connected-frequency errors. Correlation is descriptive and cannot pass
this gate.

The sequence of freezes separates method selection, discovery, and holdout.
Numerical pilots may select only predeclared resolution parameters. Direct
discovery determines which effects exceed their uncertainty. The connected
implementation then reconstructs every resolved case. Finally, action
coalitions are locked without inspecting modal consequence and evaluated on
independent holdout seeds.

\subsection{Modal Consequence and Stress Ladder}

For an eigenvalue $\lambda=\alpha+j\beta$, modal frequency and damping are
$f=|\beta|/(2\pi)$ and $\zeta=-\alpha/|\lambda|$. Left and right
eigenvectors are biorthogonally normalized. Tracking uses a registered
biorthogonal MAC (BMAC), eigenvalue distance, one-to-one assignment, and
adaptive homotopy. A modal family must be locked in discovery before the
holdout is viewed.

The finite damping externality $\ext{S}\zeta$ is computed with the same
inclusion--exclusion rule as \eqref{eq:mobius}. The registered materiality
threshold is $|\ext{S}\zeta|\ge 0.0025$ at four or more of eight holdout
points, with at least six valid tracks and same-sign fraction at least 0.75.
This 0.0025 value and the separate 5\% minimum-damping screen are study
criteria, not universal grid-code limits.

Two increasingly targeted stress tests were frozen after the nominal
materiality claim failed. E05C scales load to the first 0.85-pu voltage
boundary using eight new seeds and all 15 retained coalitions. E05D was
intended to vary grid strength, but its immutable rule required a baseline
minimum damping of at least 5\% before any coalition could be evaluated. A
later descriptive audit at $\kappa_{\mathrm{grid}}=1$ was explicitly barred
from changing any gate.

\subsection{Evidence Ladder}

Table~\ref{tab:evidence} distinguishes structural evidence from operational
evidence. This ordering prevents an exact mathematical decomposition from
being promoted automatically into a stability or mitigation claim.

\begin{table}[t]
\caption{Registered Evidence Ladder and Terminal Status}
\label{tab:evidence}
\centering
\footnotesize
\begin{tabular}{p{0.08\columnwidth}p{0.54\columnwidth}p{0.20\columnwidth}}
\toprule
ID & Question & Status\\
\midrule
C1 & Do finite re-equilibrated pair/triple externalities exist? & Supported\\
C2 & Does connected reconstruction close the direct result? & Supported\\
D1 & Does exact Schur factor anatomy close? & Supported\\
D2 & Are finite-pole externalities reproducible? & Supported\\
D3 & Can certified contours localize pole motion? & Supported, conditional\\
N1 & Is nominal modal materiality established? & Rejected\\
N2 & Is load-stress materiality established? & Rejected\\
N3 & Is weak-grid materiality established? & Unresolved; baseline ineligible\\
N4 & Is a cycle/SCC intervention established? & Unassessed\\
\bottomrule
\end{tabular}
\end{table}
